\subsection{Reporte de limpieza de datos}

Hicimos la limpieza sobre los problemas que encontramos en 2.4. Para cada uno decidimos si lo corregimos, lo eliminamos o lo ignoramos, que son las tres estrategias del curso. Ningún archivo original se borra: lo que se elimina solo deja de usarse y queda registrado en \nolinkurl{fotos.csv} (columnas \texttt{excluida} y \texttt{motivo}).

\begin{table}[H]
  \centering
  \small
  \caption{Problemas de calidad y cómo los tratamos.}
  \label{tab:limpieza}
  \begin{tabularx}{\textwidth}{>{\raggedright\arraybackslash}p{0.22\textwidth}>{\raggedright\arraybackslash}p{0.11\textwidth}>{\raggedright\arraybackslash}X>{\raggedright\arraybackslash}p{0.25\textwidth}}
    \toprule
    Problema (2.4) & Estrategia & Qué hicimos & Resultado \\
    \midrule
    Copias exactas dentro de una misma vaca (carpetas descargadas dos veces y nombres tipo \nolinkurl{IMG_0530(1).HEIC}) & Eliminar & Dejamos una sola copia por huella digital (SHA-256) dentro de cada vaca. Cuando hay varias, nos quedamos con la que no lleva ``(1)''. & De 477 archivos se quitaron 122 copias y quedaron 355 fotos únicas. \\
    \addlinespace
    Fotos idénticas en dos vacas (1548 y 2107) & Eliminar & Quitamos esas fotos de las dos vacas porque no sabemos a cuál pertenecen. & Se excluyeron 8 filas de fotos (4 fotos en 2 vacas). La 1548 conserva sus 3 fotos propias y la 2107 se queda sin fotos. \\
    \addlinespace
    Vaca 2123 con dos carpetas de fotos distintas que se enciman & Eliminar & Excluimos la vaca completa, porque no se sabe qué fotos van con cada una de sus dos evaluaciones y adivinar podría asignar un BCS equivocado. & Se excluyeron 14 fotos y 2 evaluaciones. \\
    \addlinespace
    Hora de la bitácora con ``x'' o con otro formato & Corregir e ignorar & ``13 00'' pasa a ``13:00''. Las ``x'' y los formatos raros quedan vacíos y el valor original se guarda en \texttt{hora\_original}. La hora no se usa en el modelo. & De 83 filas, 16 quedan con la hora vacía (14 con ``x'' y 2 con otro formato). \\
    \addlinespace
    BCS con coma decimal o fuera de escala & Corregir y verificar & El BCS se convierte a número y se revisa que esté entre 1 y 5 en pasos de 0.25. & Ningún BCS quedó sin convertir, fuera de rango ni fuera de los pasos de 0.25. \\
    \addlinespace
    Fotos del Samsung S22 con orientación EXIF distinta de 1 & Corregir & Los JPG con orientación distinta de 1 se enderezan. En los HEIC la librería ya aplica la rotación al abrirlos. & 20 JPG se enderezaron. Los otros 4 JPG ya estaban derechos y se copiaron sin cambios. \\
    \addlinespace
    Dos formatos (HEIC y JPG) y varias resoluciones & Corregir & Todo pasa a JPG (véase 3.5). La resolución no se cambia en esta fase. & Véase la sección~\ref{sec:reformateados}. \\
    \addlinespace
    14 vacas con dos evaluaciones de BCS distintas & Ignorar, documentado & No eliminamos ninguna evaluación. Guardamos las dos y entrenamos con el promedio (véase 3.3.1). & 13 de esas vacas siguen en el dataset, porque la 2123 se excluyó. \\
    \addlinespace
    Valores faltantes & Ignorar, documentado & No imputamos nada, porque no tenemos campos numéricos con faltantes que estimar. & No aplica. \\
    \bottomrule
  \end{tabularx}
\end{table}

Cifras de la limpieza:

\begin{itemize}
  \item \textbf{Fotos:} se leyeron 477 archivos, se quitaron 122 copias exactas y quedaron 355 fotos únicas. Se excluyeron 22 (14 de la vaca 2123 y 8 filas repetidas entre la 1548 y la 2107), así que siguen 333 fotos de 67 vacas.
  \item \textbf{Bitácora:} tiene 83 evaluaciones, se excluyeron 2 (las de la 2123) y la bitácora limpia queda con 81 evaluaciones.
\end{itemize}

\subsubsection*{Pendientes que siguen abiertos y su posible efecto}

\begin{itemize}
  \item Las 7 vacas con etiqueta ``fresca'' contradictoria entre los expertos y el conjunto de datos de VMS 3 siguen sin auditarse con el CAETEC. Si alguna etiqueta está mal, afecta directamente al BCS de entrenamiento y a las métricas.
  \item Al excluir la 2123 y la 2107 perdemos vacas y fotos en un dataset que ya es chico. \textbf{Pendiente:} confirmar en equipo la exclusión de la vaca 2123.
  \item Hay clases de BCS con muy pocas vacas (2.25 tiene una sola evaluación), así que las métricas pueden moverse mucho con un solo error.
  \item Las vistas de las 333 fotos las marcó una sola persona del equipo el 26/10/2026. \textbf{Pendiente:} definir si otra persona revisa una muestra de las marcas.
\end{itemize}
